\documentclass[]{tufte-handout}

% ams
\usepackage{amssymb,amsmath}

\usepackage{ifxetex,ifluatex}
\usepackage{fixltx2e} % provides \textsubscript
\ifnum 0\ifxetex 1\fi\ifluatex 1\fi=0 % if pdftex
  \usepackage[T1]{fontenc}
  \usepackage[utf8]{inputenc}
\else % if luatex or xelatex
  \makeatletter
  \@ifpackageloaded{fontspec}{}{\usepackage{fontspec}}
  \makeatother
  \defaultfontfeatures{Ligatures=TeX,Scale=MatchLowercase}
  \makeatletter
  \@ifpackageloaded{soul}{
     \renewcommand\allcapsspacing[1]{{\addfontfeature{LetterSpace=15}#1}}
     \renewcommand\smallcapsspacing[1]{{\addfontfeature{LetterSpace=10}#1}}
   }{}
  \makeatother

\fi

% graphix
\usepackage{graphicx}
\setkeys{Gin}{width=\linewidth,totalheight=\textheight,keepaspectratio}

% booktabs
\usepackage{booktabs}

% url
\usepackage{url}

% hyperref
\usepackage{hyperref}

% units.
\usepackage{units}


\setcounter{secnumdepth}{-1}

% citations

% pandoc syntax highlighting

% longtable

% multiplecol
\usepackage{multicol}

% strikeout
\usepackage[normalem]{ulem}

% morefloats
\usepackage{morefloats}


% tightlist macro required by pandoc >= 1.14
\providecommand{\tightlist}{%
  \setlength{\itemsep}{0pt}\setlength{\parskip}{0pt}}

% title / author / date
\date{}


\begin{document}





\pagenumbering{gobble} \vspace*{30pt}

\huge FOREST MANAGEMENT PLAN

\vspace{30pt}

\LARGE Two Dog Timber Property

\vspace{22pt}

\normalsize Town of Harrietstown, Franklin County, New York

\vspace{70pt}

\small 157 acres

\vspace{18pt}

Parcel IDs:

\vspace{8pt}

457.-4-17.100

\vspace{10pt}

Deed (Liber/Pages): 725/254

\vspace{10pt}

480-a certification number: 16-062

\vspace{10pt}

Original certication date: 12/08/2014

\vspace{60pt}

Prepared by:

\includegraphics[width=0.5\linewidth,height=0.5\textheight]{C:/Users/Neal/projects/cruise/logo-name-small}

John D. Foppert \& Neal F. Maker\\
96 Durocher Road\\
Saranac, NY 12981

\vspace{30pt}

\large February 18, 2020

\pagebreak
\pagenumbering{arabic}

\section{Introduction}\label{introduction}

This plan covers the fifteen year period from 2019 to 2034. It lays out
the near- and medium-term actions that should guide the development of
the Two Dog Timber forest. It also qualifies portions of the property
for continued enrollment in the 480-a Forest Tax Law program and
commensurate reduction in property taxes.\footnote{Further information
  about 480-a can be found at the New York Department of Conservation's
  website: \url{https://www.dec.ny.gov/lands/5236.html}. \vspace{20pt}}
Owners participating in the program are obliged to manage enrolled
portions of their property according to their approved forest management
plan and to make any reasonable investments for improvement that the
plan recommends. Its recommendations were developed in accordance with
the principles and practices of scientifically sound forestry, as
described in the relevant management guidelines, textbooks and academic
journals.

\section{Property Description}\label{property-description}

The Two Dog Timber forest represents a well-managed resource positioned
to continue generating value on a sustainable basis. Ninety seven
percent of the 157 acre property is productive forestland to be managed
according to this plan. The property is located in the Adirondack Park
in the Town of Harrietstown, Franklin County, New York. The property is
accessed off of Kiwassa Lake Road and State Route 3. The productive
forestland is operationally accessible via a well developed internal
road network. Property-wide, elevations range from 1570 to 1870 feet
above mean sea level. The property includes some areas of frequently wet
soils, but there are no mapped wetlands or streams on the property. The
property's boundaries are well established. To the east, the property
fronts Kiwassa Lake Road. It is bounded elsewhere by property lines
running just over two miles. These lines are generally well marked in
the field (typically with tree blazes and red boundary marking paint),
but will require on-going maintenance. Soils, forest health, and other
pertinent topics are discussed in the individual stand area descriptions
that follow.

\section{Principles, Goals \& Strategies For Forest
Management}\label{principles-goals-strategies-for-forest-management}

\subsection{Conservation}\label{conservation}

The ecological functioning, productive capacity and biological diversity
of the forest resource should be maintained or improved over time so as
to provide opportunities for the current or future landowners to
continue to enjoy and use the property. A management strategy that is
sustainable in the long-term and viable in the short- and medium-terms
offers a strong measure of protection against future development or
conversion.

\subsection{Timber management}\label{timber-management}

Management should provide regular returns from timber harvesting.
Long-term value growth is provided by maintaining full site occupancy
with investment-grade stems: healthy trees capable of producing high
quality sawtimber or veneer and worth retaining in the stand until they
reach their full, site- and species-specific target diameters. Tree
species which yield sought-after, high-value wood should be promoted
within each stand or, when regenerating a new stand, attention should be
paid to providing the stand conditions which favor the establishment of
those species. At a property-wide scale, a variety of species should be
maintained, providing options for seizing future market opportunities
and a hedge against species-specific market depreciation. Among desired
species, additional preference should be given to individual trees of
sufficient vigor and grade-potential for strong future value growth.
Consideration of economic efficiency should inform the timing and
coordination of infrastructure investments and stand maintenance,
improvement and harvest operations.

\subsection{Amenity values}\label{amenity-values}

Conscientious management can create or maintain a landscape that is
attractive, accessible and conducive to reflection, exploration and
appreciation. Attractiveness can be managed for by fostering diversity
within the landscape: promoting the growth and development of the most
appealing individual trees in some places; maintaining the look, feel
and accompanying privacy provided by a dense forest in other places; and
elsewhere creating occasional vistas out from the forest and
improvements in depth of visual penetration within it. Carefully planned
and deliberately located infrastructure should facilitate the satisfying
use of the property, creating an appropriate balance between access and
connectedness, on the one hand, and places of refuge and sanctuary, on
the other. A system of roads and trails of various sizes, suited for
various purposes, and interconnected with a broader trail network,
provide for both enjoyable recreation and efficient operations.

\section{Stand Descriptions \& Management
Recommendations}\label{stand-descriptions-management-recommendations}

Presented below are detailed stand-by-stand descriptions of the forest,
the long-term structural, compositional and functional goals for each
stand, and the near-term silvicultural treatments or management
activities that have been prescribed to advance each stand toward those
goals. The data presented in the following pages was obtained from a
thorough inventory of the property in October of 2019. General
conditions were assessed qualitatively in conjunction with quantitative
sampling. Observational notes and sample summary statistics together
provide the basis for the stand descriptions and management
recommendations. All sampling was done using a systematic sample and
variable radius plots. In stands with uneven-aged structures, all trees
6" in diameter at breast height (dbh) and larger were measured in each
plot. In stands with even-aged structures, all main-canopy trees were
measured in each plot.

When contractors are used to implement silvicultural prescriptions, they
should be highly skilled, properly equipped, fully insured, and closely
supervised. A professional forester should prepare and administer
commercial treatments, and logging operations should be timed to
coincide with favorable weather conditions (working on wet soils only
when they are frozen, for instance) and favorable timber markets. The
dates assigned to timber harvests and other management activities
prescribed in this plan are intended to guide, rather than constrain,
forest management. To accommodate dynamic markets and variable weather,
scheduled timber harvests may be advanced or delayed by one year from
the date indentified in this plan; if operational or economic conditions
change substantially, the management schedule may be further revised by
an ammendment to this plan.

The property should be reassessed in 2024 and the findings brought to
bear on a reassessment of the goals and strategies proposed in this
plan, leading to a formal management plan update.

\newpage

\section{Management Schedule}\label{management-schedule}

\small 2020 - Management plan update; boundary line maintenance

\vspace{5pt}

\noindent 2021 - No scheduled activity

\vspace{5pt}

\noindent 2022 - No scheduled activity

\vspace{5pt}

\noindent 2023 - No scheduled activity

\vspace{5pt}

\noindent 2024 - No scheduled activity

\vspace{5pt}

\noindent 2025 - Management plan update; boundary line maintenance

\vspace{5pt}

\noindent 2026 - No scheduled activity

\vspace{5pt}

\noindent 2027 - No scheduled activity

\vspace{5pt}

\noindent 2028 - No scheduled activity

\vspace{5pt}

\noindent 2029 - No scheduled activity

\vspace{5pt}

\noindent 2030 - Management plan update; boundary line maintenance

\vspace{5pt}

\noindent 2031 - No scheduled activity

\vspace{5pt}

\noindent 2032 - No scheduled activity

\vspace{5pt}

\noindent 2033 - No scheduled activity

\vspace{5pt}

\noindent 2034 - Full management plan revision; boundary line
maintenance

\newpage

\section{Area 1}\label{area-1}

Mixedwood\\
\noindent 152.80 acres total\\
\noindent 0.00 acres ineligible wetlands\\
\noindent 1.20 acres ineligible roads and landings

\subsection{Site-specific information}\label{site-specific-information}

\begin{quote}
\begin{itemize}
\tightlist
\item
  \textbf{Soils:}\\
  \indent\indent Monadnock, Tunbridge, Becket
\end{itemize}
\end{quote}

\begin{quote}
\begin{itemize}
\tightlist
\item
  \textbf{Site Class:}\\
  \vspace{2pt} II (determined from soil mapping and field assessment)
\end{itemize}
\end{quote}

\begin{quote}
\begin{itemize}
\tightlist
\item
  \textbf{Access:}\\
  \vspace{2pt} Less than 1 mile
\end{itemize}
\end{quote}

\begin{quote}
\begin{itemize}
\tightlist
\item
  \textbf{Stand history:}\\
  \vspace{2pt} This stand has been continuously forested for at least
  the past 100 years. It appears to have been harvested on a recurring
  basis over most of that period, including a cut of mostly spruce in
  around the turn of the last century, maple and birch in the 1950's, a
  hemlock-focused harvest in the 1990's. A silviculturlly prescribed
  harvest was initiated in 2016 and 2017, remoing 46.635 MBF of hardwood
  sawlogs and 60 cords of firewood. The harvest was completed in 2018,
  removing an additional 21.755 MBF of hardwood logs, 7.310 MBF of
  hemlock logs, and 117.4 cords of firewood and pulp.
\end{itemize}
\end{quote}

\subsection{Current forest
information}\label{current-forest-information}

\begin{quote}
\begin{itemize}
\tightlist
\item
  \textbf{Age Class Structure:}\\
  \vspace{2pt} Uneven-aged\\

  \begin{marginfigure}
  \includegraphics{fmp-ny_files/figure-latex/unnamed-chunk-2-1} \caption[Distributions are approximated with kernel density estimation]{Distributions are approximated with kernel density estimation. Common species are those that account for at least 8 percent of the total stocking and areas under each curve represent species basal areas.}\label{fig:unnamed-chunk-2}
  \end{marginfigure}
\end{itemize}
\end{quote}

\begin{quote}
\begin{itemize}
\tightlist
\item
  \textbf{Species (\% stocking):}\\
  \vspace{2pt} hard maple (34\%), hemlock (22\%), yellow birch (10\%),
  beech (10\%), soft maple (6\%), black cherry (5\%), ash (4\%),
  basswood (3\%), spruce (3\%), fir (2\%), hophornbeam (1\%), paper
  birch (1\%), cedar (0\%)
\end{itemize}
\end{quote}

\begin{quote}
\begin{itemize}
\tightlist
\item
  \textbf{Regeneration:}\\
  \vspace{2pt} Regeneration is present across approximately 2/3 of the
  stand. Over half of this area, it is concentrated among beech
  saplings. In places, the high density of beech regeneration excludes
  the establishment of any other species. Nevertheless, approximately
  1/3 of plots samples included desirable hardwood species positioned to
  survive into the future stand.
\end{itemize}
\end{quote}

\begin{quote}
\begin{itemize}
\tightlist
\item
  \textbf{Forest health:}\\
  \vspace{2pt} Few forest health issues were observed on this portion of
  the property. Beech bark disease is present and affects older beech
  trees. Windthrow will be a persistent concern in portions of the stand
  where soils are wet or shallow-to-bedrock.
\end{itemize}
\end{quote}

\begin{quote}
\begin{itemize}
\tightlist
\item
  \textbf{Size class structure (\%BA):}\\
  \vspace{2pt} \indent 6-10'': 24\% \textbar{} 11-16'': 45\% \textbar{}
  17-22'': 26\% \textbar{} 23"+: 5\%
\end{itemize}
\end{quote}

\subsection{Inventory information}\label{inventory-information}

\begin{quote}
\begin{itemize}
\tightlist
\item
  32 points, 10 BAF, October, 2019
\end{itemize}
\end{quote}

\begin{figure}
\includegraphics{fmp-ny_files/figure-latex/unnamed-chunk-5-1} \caption[Points represent individual plots]{Points represent individual plots. Asterisk represnts stand average. Radial lines are quadratic stand diameters.}\label{fig:unnamed-chunk-5}
\end{figure}

\begin{table}

\caption{\label{tab:unnamed-chunk-6}Measures of stocking for all live trees (total) and acceptable 
      growing stock.}
\centering
\begin{tabular}[t]{lrr}
\toprule
  & Total & Acceptable\\
\midrule
Basal area (sqft/ac) & 113 & 98\\
QSD (in) & 10 & 10\\
Stems/ac & 195 & 164\\
\bottomrule
\end{tabular}
\end{table}

\subsection{Silvicultural
prescription}\label{silvicultural-prescription}

No silvicultural treatment is called for at this time.

\newpage



\end{document}
